% !TEX root = ../main.tex
\documentclass[../main.tex]{subfiles}

\begin{document}

\intermezzo{RoleSim: Agent-Based Modeling of Role Constellation Dynamics}

\begin{center}
\begin{tikzpicture}
    \fill[Marigold] (-2,0) rectangle (-1.7,0.08);
    \fill[Marigold!80] (-1.5,0) rectangle (-1.2,0.08);
    \fill[Marigold!60] (-1,0) rectangle (-0.7,0.08);
    \fill[Marigold!40] (-0.5,0) rectangle (-0.2,0.08);
    \fill[Marigold!20] (0,0) rectangle (0.3,0.08);
\end{tikzpicture}
\end{center}

\vspace{0.5cm}

\begin{chapteroverview}
This chapter introduces \RoleSim, an agent-based model that complements empirical cartography with counterfactual exploration. Where \RoleBox maps observed constellations, \RoleSim asks ``what if?''---exploring how constellations might evolve under different conditions, testing mechanism sufficiency, and projecting possible futures.
\end{chapteroverview}


%═══════════════════════════════════════════════════════════════════════════════
\section*{Why Simulation?}
%═══════════════════════════════════════════════════════════════════════════════

Empirical observation reveals where actors \emph{are}. Simulation explores where they \emph{could be}:

\sidenote{\textbf{Counterfactual Questions:} What if a key intermediary exited? What if policy conditions shifted? What if fitness landscapes tilted?}

\begin{description}[leftmargin=0.5cm, itemsep=0.3em]
\item[Counterfactual Exploration] What trajectories were possible but not taken?
\item[Mechanism Testing] Do CAS mechanisms generate observed patterns?
\item[Sensitivity Analysis] Which parameters most strongly shape outcomes?
\item[Future Projection] Given current configurations, what trajectories are likely?
\end{description}


%═══════════════════════════════════════════════════════════════════════════════
\section*{Model Architecture}
%═══════════════════════════════════════════════════════════════════════════════

\RoleSim instantiates CAS mechanisms computationally:

\subsection*{Agents}

Each agent represents an intermediary organization with:
\begin{itemize}[leftmargin=*, itemsep=0.2em]
\item \textbf{Role profile}: Multi-dimensional vector in role-space
\item \textbf{Resources}: Capacity for activities and adaptation
\item \textbf{Network position}: Ties to other agents
\item \textbf{Adaptive behavior}: Rules for adjusting profile based on feedback
\end{itemize}

\subsection*{Fitness Landscape}

The fitness function determines agent viability:
\begin{itemize}[leftmargin=*, itemsep=0.2em]
\item Environmental conditions (policy, technology)
\item Niche overlap with other agents
\item Network position value
\item Resource requirements
\end{itemize}

\subsection*{Dynamics}

Each timestep:
\begin{enumerate}[leftmargin=*, itemsep=0.1em]
\item Agents evaluate fitness at current position
\item Agents observe neighbors' success
\item Agents adapt profiles toward higher-fitness regions
\item Network ties form/dissolve based on similarity and complementarity
\item Environment may shift (exogenous shocks)
\end{enumerate}


%═══════════════════════════════════════════════════════════════════════════════
\section*{Parameterization from RoleBox}
%═══════════════════════════════════════════════════════════════════════════════

\RoleSim is grounded in empirical data:

\begin{itemize}[leftmargin=*, itemsep=0.2em]
\item Initial agent configurations from observed EIT ecosystem
\item Fitness function shapes derived from survival/growth patterns
\item Network structures from measured relationships
\item Environmental parameters from policy and technology trajectories
\end{itemize}

\sidenote{\textbf{Grounding:} Simulation explores variations on observed reality, not arbitrary parameter spaces.}


%═══════════════════════════════════════════════════════════════════════════════
\section*{Analytical Uses}
%═══════════════════════════════════════════════════════════════════════════════

\begin{description}[leftmargin=0.5cm, itemsep=0.3em]
\item[Knockout Experiments] Remove successful agents; observe restructuring
\item[Policy Scenarios] Shift environmental parameters; trace constellation response
\item[Mechanism Isolation] Enable/disable specific CAS mechanisms; assess contribution
\item[Parameter Sweeps] Systematic variation to identify sensitive parameters
\end{description}


%═══════════════════════════════════════════════════════════════════════════════
\section*{Chapter Summary}
%═══════════════════════════════════════════════════════════════════════════════

\begin{summarybox}
\begin{itemize}[leftmargin=*, itemsep=0.3em]
\item \textbf{\RoleSim} complements empirical cartography with counterfactual exploration.

\item \textbf{Agent-based modeling} instantiates CAS mechanisms computationally.

\item \textbf{Parameterization} from \RoleBox data grounds simulation in observed reality.

\item \textbf{Analytical uses} include knockout experiments, policy scenarios, and mechanism testing.

\item \textbf{Together}, \RoleBox and \RoleSim enable understanding of both actual and possible constellation dynamics.
\end{itemize}
\end{summarybox}

\vspace{0.5cm}

\begin{center}
\begin{tikzpicture}
    \fill[Marigold] (-2,0) rectangle (-1.7,0.08);
    \fill[Marigold!80] (-1.5,0) rectangle (-1.2,0.08);
    \fill[Marigold!60] (-1,0) rectangle (-0.7,0.08);
    \fill[Marigold!40] (-0.5,0) rectangle (-0.2,0.08);
    \fill[Marigold!20] (0,0) rectangle (0.3,0.08);
\end{tikzpicture}
\end{center}

\closeintermezzo

\end{document}
